\section*{Chapitre \ref{ch:b1:diagenesis-basins} --- La diagenèse et les bassins sédimentaires}

\begin{solution}{exo:b1:diagenesis-basins:1}
Grès $0.49\,e^{-2000/3704} = 0.49\times0.583 = \num{0.286}$~; argilite
$0.63\,e^{-2000/1960} = 0.63\times0.360 = \num{0.227}$. La porosité est divisée par deux quand
$e^{-z/\lambda} = 1/2$, à $z = \lambda\ln 2$~: \qty{2567}{m} pour le grès,
\qty{1359}{m} pour l'argilite.
\end{solution}

\begin{solution}{exo:b1:diagenesis-basins:2}
$\Delta A = \Delta E + S = 30 + 50 = \qty{80}{m}$. Sur cela, \qty{60}{m} sont comblés
de sédiment, donc $\Delta W = 80 - 60 = \qty{20}{m}$~: la plate-forme s'approfondit de
\qty{20}{m}.
\end{solution}

\begin{solution}{exo:b1:diagenesis-basins:3}
$h = Y(\rho_m - \rho_w)/(\rho_m - \rho_s)$~: avec \qty{2300}{kg.m^{-3}},
$1.0\times2270/1000 = \qty{2.27}{km}$~; avec \qty{2500}{kg.m^{-3}},
$1.0\times2270/800 = \qty{2.84}{km}$. Chaque mètre de sédiment prend la place d'un
mètre de \omterm{def:b1:layered-earth:layers}{manteau} dans la colonne~; plus le sédiment est dense, plus la masse
qui lui manque par rapport au \omterm{def:b1:layered-earth:layers}{manteau} est petite, et plus il en faut pour combler
le même déficit que le trou rempli d'eau.
\end{solution}

\begin{solution}{exo:b1:diagenesis-basins:4}
(a)~Calcite $100.09/2.71 = \qty{36.93}{cm^3.mol^{-1}}$~; dolomite $184.40/2.84
= \qty{64.93}{cm^3.mol^{-1}}$. (b)~Deux moles de calcite occupent
\qty{73.86}{cm^3}, une de dolomite \qty{64.93}{cm^3}~: le solide rétrécit de
$1 - 64.93/73.86 = \qty{12.1}{\%}$. (c)~Dans un mètre cube, les \qty{0.90}{m^3}
de calcite deviennent $0.90\times(1 - 0.121) = \qty{0.791}{m^3}$ de dolomite~; la
porosité devient $1 - 0.791 = \num{0.21}$.
\end{solution}

\begin{solution}{exo:b1:diagenesis-basins:5}
(a)~$s = 1000 - 0.63\times1960\times(1 - e^{-1000/1960}) = \qty{506.5}{m}$~; la
porosité moyenne est $1 - 506.5/1000 = \num{0.493}$. (b)~Le toit étant à
\qty{2000}{m}, la résolution de $s(2000, b) = 506.5$ donne $b = \qty{2628.7}{m}$~: la
couche a \qty{628.7}{m} d'épaisseur, de porosité moyenne $1 - 506.5/628.7 = \num{0.194}$.
(c)~Les grains sont les mêmes, donc l'eau perdue est l'épaisseur perdue~:
$1000 - 628.7 = \qty{371}{m^3}$ d'eau par mètre carré.
\end{solution}

\begin{solution}{exo:b1:diagenesis-basins:6}
(a)~$s = 300 - 0.49\times3704\times(e^{-2500/3704} - e^{-2800/3704}) =
\qty{228.10}{m}$. (b)~À partir de $b_0 = 300$~: $s(0, 300) = 158.80$, $1 - \varphi(300) =
0.5481$, $b_1 = 300 + (228.10 - 158.80)/0.5481 = \qty{426.44}{m}$~; puis
$s(0, 426.44) = 229.07$, $1 - \varphi = 0.5633$, $b_2 = 426.44 - 0.96/0.5633 =
\qty{424.73}{m}$~; un troisième pas la modifie de moins d'un millimètre. Le
grès avait \qty{424.7}{m} d'épaisseur à la surface. (c)~$424.7/300 = 1.42$~: l'enfouissement
l'a aminci d'un facteur \num{1.42}.
\end{solution}

\begin{solution}{exo:b1:diagenesis-basins:7}
(a)~La glace contient $52\times10^6\times916.7/1000 = \qty{4.77e7}{km^3}$ d'eau~;
étalée sur \qty{361.84e6}{km^2} cela fait $\qty{0.1317}{km} = \qty{132}{m}$ de niveau
de la mer. (b)~Les rivages donnent \qty{134}{m}, soit \qty{2}{m} de plus. (c)~$4.767\times
10^7/0.134 = \qty{355.7e6}{km^2}$, soit \qty{1.7}{\%} de moins qu'aujourd'hui. Quand le niveau de la mer
a baissé, les plateaux continentaux ont émergé, si bien que l'océan glaciaire était plus petit, et le
même volume d'eau qui lui était prélevé l'abaissait davantage.
\end{solution}

\begin{solution}{exo:b1:diagenesis-basins:8}
(a)~$(-9.1 - (-98.1))/8 = \qty{11.1}{m}$ par millier d'années, proche des
\qty{12}{m} trouvés d'après les rivages. (b)~\qty{40}{mm} par an font \qty{40}{m} par
millier d'années, soit $40/12 = 3.3$ fois la moyenne~; en cinq cents ans,
$0.040\times500 = \qty{20}{m}$.
\end{solution}

\begin{solution}{exo:b1:diagenesis-basins:9}
(a)~$S = 50\times2 = \qty{100}{m}$ et $\Delta E = \qty{-40}{m}$~:
$\Delta A = \qty{60}{m}$, entièrement comblés de sédiment puisque la profondeur est inchangée.
(b)~En \num{20000} ans le socle ne s'enfonce que de $50\times0.02 = \qty{1}{m}$~:
le \omterm{def:b1:diagenesis-basins:eustasy}{niveau marin relatif} baisse de $120 - 1 = \qty{119}{m}$. Le fond marin, qui se trouvait
\qty{20}{m} sous le niveau de la mer, se retrouve $119 - 20 = \qty{99}{m}$ au-dessus de lui~: toute la
plate-forme émerge.
\end{solution}

\begin{solution}{exo:b1:diagenesis-basins:10}
(a)~Pour une couche mince la porosité est presque constante sur son épaisseur~: son épaisseur
solide est $s \approx h(1 - \varphi(z))$ en profondeur et $s \approx h_0(1 -
\varphi_0)$ à la surface. En égalant les deux on obtient $h_0 \approx h(1 -
\varphi(z))/(1 - \varphi_0)$. (b)~$\varphi(3000) = 0.63\,e^{-3000/1960} =
0.136$, donc $h_0 \approx 10\times0.864/0.37 = \qty{23.3}{m}$, contre
\qty{23.1}{m} exactement. (c)~Ramenée à la surface, la couche a
\qty{23}{m} d'épaisseur, elle n'est plus mince devant la décroissance de la porosité près de la
surface~: sa porosité passe de $\varphi_0$ à son toit à moins à sa base, si bien que
ses grains tiennent dans un peu moins que l'épaisseur calculée avec $\varphi_0$
partout.
\end{solution}

\begin{solution}{exo:b1:diagenesis-basins:11}
(a)~$s = 3000 - 0.63\times1960\times(1 - e^{-3000/1960}) = \qty{2032}{m}$~; la
porosité moyenne est $1 - 2032/3000 = \num{0.323}$ et la masse volumique moyenne
$0.323\times1030 + 0.677\times2700 = \qty{2161}{kg.m^{-3}}$. (b)~Le facteur est
$(3300 - 1030)/(3300 - 2161) = \num{1.99}$~; le remplissage représente une \omterm{def:b1:diagenesis-basins:subsidence}{subsidence}
sous charge d'eau de $3000/1.99 = \qty{1505}{m}$. (c)~Avec la masse volumique des grains le facteur
serait $2270/600 = \num{3.78}$ et la \omterm{def:b1:diagenesis-basins:subsidence}{subsidence} seulement \qty{793}{m}. Ce serait
faux parce qu'un tiers du remplissage est de l'eau~: le sédiment est bien plus léger que ses
grains, et seule la masse volumique du sédiment saturé compte dans le bilan des
masses.
\end{solution}

\begin{solution}{exo:b1:diagenesis-basins:12}
(a)~Le cube contient \qty{0.90}{m^3} $= \qty{9.0e5}{cm^3}$ de calcite, soit
$9.0\times10^5/36.93 = \num{24370}$ moles. (b)~Un magnésium pour deux calcites~:
\num{12180} moles. (c)~Un kilogramme d'eau de mer contient $35.165\times0.0365 =
\qty{1.284}{g}$ de magnésium, $1.284/24.305 = \qty{0.0528}{mol}$~; la réaction
exige $12180/0.0528 = \qty{2.31e5}{kg}$ d'eau de mer, $\qty{224}{m^3}$ --- plus
de deux cents fois le volume de la roche, même si chaque ion magnésium est
capté. (d)~La dolomite ne peut se former à grande échelle que là où de très grands volumes
d'eau traversent la roche pendant longtemps, ou là où l'eau est une saumure
plus riche en magnésium que l'eau de mer~: dans des plates-formes balayées par des saumures qui s'enfoncent,
par exemple, plutôt que dans un sédiment qui garde la même eau interstitielle.
\end{solution}

\begin{solution}{pb:b1:diagenesis-basins:1}
\textbf{1.} Grès~: $\varphi(0) = \num{0.490}$, $\varphi(800) =
0.49\,e^{-800/3704} = \num{0.395}$. Argilite~: $\varphi(800) = \num{0.419}$,
$\varphi(2600) = \num{0.167}$. Craie~: $\varphi(2600) = \num{0.110}$,
$\varphi(3200) = \num{0.072}$.

\textbf{2.} Avec $s$ de la \cref{prop:b1:diagenesis-basins:solid} et $\bar\varphi
= 1 - s/h$ (question~7)~: grès \num{0.441}, argilite \num{0.274}, craie
\num{0.090}.

\textbf{3.} $h\bar\varphi$~: grès \qty{353}{m^3}, argilite \qty{493}{m^3}, craie
\qty{54}{m^3} d'eau par mètre carré.

\textbf{4.} $\bar\varphi\rho_w + (1 - \bar\varphi)\rho_g$~: grès
$0.441\times1030 + 0.559\times2650 = \qty{1936}{kg.m^{-3}}$, argilite
\qty{2242}{kg.m^{-3}}, craie \qty{2559}{kg.m^{-3}}. La colonne pèse
$1936\times800 + 2242\times1800 + 2559\times600 = \qty{7.12e6}{kg}$ par mètre
carré~: une masse volumique moyenne de \qty{2225}{kg.m^{-3}}.

\textbf{5.} Lithostatique $2225\times9.81\times3200 = \qty{69.9}{MPa}$~; hydrostatique
$1030\times9.81\times3200 = \qty{32.3}{MPa}$~; \omterm{def:b1:stress-strain-rheology:lithostatic}{contrainte effective}
$69.9 - 32.3 = \qty{37.5}{MPa}$.

\textbf{6.} $0.49\,e^{-z/3704} = 0.63\,e^{-z/1960}$ donne $z = \ln(0.63/0.49)
/(1/1960 - 1/3704) = \qty{1046}{m}$. En dessous, le grès est le plus poreux~:
l'argilite commence plus humide mais perd son eau sur une distance plus courte.

\textbf{7.} Grès $800 - 0.49\times3704\times(1 - e^{-800/3704}) =
\qty{447.4}{m}$~; argilite \qty{1306.7}{m}~; craie \qty{546.0}{m}.

\textbf{8.} \qty{2300.2}{m} de grains~; l'eau fait $1 - 2300.2/3200 =
\qty{28}{\%}$ de la colonne.

\textbf{9.} Grès $447.4/15 = \qty{29.8}{m}$, argilite $1306.7/40 =
\qty{32.7}{m}$, craie $546.0/20 = \qty{27.3}{m}$ de solide par million d'années.

\textbf{10.} L'argilite a eu l'apport le plus rapide, la craie le plus lent --- mais les
trois ne diffèrent que d'un cinquième environ, bien moins que ne le laissent penser leurs
épaisseurs actuelles.

\textbf{11.} \qty{800}{m}~: son toit est à la surface aujourd'hui, là où son dépôt
vient de s'achever, donc son épaisseur au dépôt est son épaisseur actuelle.

\textbf{12.} La résolution de $s(0, b) = 1306.7$ pour l'argilite par la
\cref{met:b1:diagenesis-basins:decompaction} donne $b = \qty{2124}{m}$~: l'argilite
avait \qty{2124}{m} d'épaisseur à \qty{15}{\Ma}.

\textbf{13.} $s(0, b) = 546.0$ pour la craie donne \qty{1071}{m} à
\qty{55}{\Ma}.

\textbf{14.} Avec son toit à \qty{2124}{m}, $s(2124, b) = 546.0$ donne $b =
\qty{2748}{m}$~: la craie avait \qty{624}{m} d'épaisseur.

\textbf{15.} \qty{1071}{m} à \qty{55}{\Ma}~; $2124 + 624 = \qty{2748}{m}$ à
\qty{15}{\Ma}.

\textbf{16.} D'après les épaisseurs actuelles~: grès $800/15 = \qty{53.3}{m}$,
argilite $1800/40 = \qty{45.0}{m}$, craie $600/20 = \qty{30.0}{m}$ par million
d'années. D'après les épaisseurs au dépôt~: \qty{53.3}{m}, $2124/40 =
\qty{53.1}{m}$ et $1071/20 = \qty{53.6}{m}$. Le second jeu est le plus significatif
--- la vitesse à laquelle le fond marin a été construit --- et il montre une vitesse
régulière que les épaisseurs actuelles cachent.

\textbf{17.} $1 - 1800/2124 = \qty{15}{\%}$.

\textbf{18.} La profondeur d'eau reste la même au début et à la fin de chaque
intervalle~; l'\omterm{def:b1:diagenesis-basins:accommodation}{accommodation} créée vaut la croissance de la colonne au-dessus du
repère~: \qty{1071}{m} de 75 à \qty{55}{\Ma}, $2748 - 1071 = \qty{1677}{m}$ de
55 à \qty{15}{\Ma}, et $3200 - 2748 = \qty{452}{m}$ depuis \qty{15}{\Ma}.

\textbf{19.} $S = \Delta A - \Delta E$~: $1677 - (-60) = \qty{1737}{m}$ de 55 à
\qty{15}{\Ma}~; $452 - (-30) = \qty{482}{m}$ depuis.

\textbf{20.} $1737/40 = \qty{43}{m}$ et $482/15 = \qty{32}{m}$ par million
d'années.

\textbf{21.} $60/1737 = \qty{3.5}{\%}$ et $30/482 = \qty{6.2}{\%}$. Le niveau de la mer a baissé
dans les deux intervalles, donc l'\omterm{def:b1:diagenesis-basins:eustasy}{eustatisme} a gêné le comblement~: le socle a dû s'enfoncer
d'autant plus pour faire la même place. La \omterm{def:b1:diagenesis-basins:subsidence}{subsidence} a fait pratiquement toute
l'\omterm{def:b1:diagenesis-basins:accommodation}{accommodation}.

\textbf{22.} L'argilite et la craie avaient \qty{2748}{m} d'épaisseur à \qty{15}{\Ma}
et en ont \qty{2400}{m} aujourd'hui~: leur \omterm{def:b1:diagenesis-basins:compaction}{compaction} a fait \qty{348}{m} de place,
soit \qty{77}{\%} des \qty{452}{m} comblés par le grès.

\textbf{23.} Chaque unité a été amincie par l'enfouissement qui l'a suivie, d'une
quantité qui dépend de sa lithologie et de la profondeur atteinte~; la place faite dans
chaque intervalle doit se mesurer avec les épaisseurs que les unités avaient alors, ce que
seule la \omterm{def:b1:diagenesis-basins:compaction}{décompaction} donne. Et même alors la \omterm{def:b1:diagenesis-basins:subsidence}{subsidence} ainsi trouvée comprend
l'enfoncement causé par le poids du sédiment, qu'il faut retirer pour atteindre la
\omterm{def:b1:diagenesis-basins:subsidence}{subsidence tectonique}.

\textbf{24.} \qty{2124}{m}, environ \qty{2.1}{km} --- \qty{324}{m} de plus que les
\qty{1800}{m} que le forage trouve aujourd'hui.
\end{solution}
